\heading{数学物理方法（下）-第1次作业}


\begin{question}
在弦的横振动问题中，若弦受到一个与速度成正比（比例系数为 \(d\)）的阻尼，
试导出弦的有阻尼振动方程。又若除了阻尼力之外，弦还受到与弦的位移成正比
（比例系数为 \(k\)）的回复力，则此时弦的振动满足的方程是什么？
\begin{solution}
设弦的线密度为 \(\rho\)，张力为 \(T\)，横位移为 \(u(x,t)\)。小振幅近似下，张力在横向的合力为
\[
T\frac{\partial^2u}{\partial x^2}\,dx.
\]
阻尼力与速度反向，单位长度阻尼力为
\[
-d\frac{\partial u}{\partial t}.
\]
由牛顿第二定律得
\[
\rho\,dx\,\frac{\partial^2u}{\partial t^2}
=T\frac{\partial^2u}{\partial x^2}\,dx
-d\frac{\partial u}{\partial t}\,dx.
\]
因此有阻尼弦振动方程为
\[
\frac{\partial^2u}{\partial t^2}
=a^2\frac{\partial^2u}{\partial x^2}
-\gamma\frac{\partial u}{\partial t},
\qquad
a^2=\frac{T}{\rho},\quad \gamma=\frac{d}{\rho}.
\]
若再有与位移成正比且方向相反的回复力 \(-ku\)，则
\[
\rho\frac{\partial^2u}{\partial t^2}
=T\frac{\partial^2u}{\partial x^2}
-d\frac{\partial u}{\partial t}-ku,
\]
即
\[
\frac{\partial^2u}{\partial t^2}
=a^2\frac{\partial^2u}{\partial x^2}
-\gamma\frac{\partial u}{\partial t}
-\omega_0^2u,
\qquad
\omega_0^2=\frac{k}{\rho}.
\]
\end{solution}
\end{question}

\begin{question}
一截面不均匀的细杆（如图），杆的截面积与离杆尖距离的平方成正比，
比例系数为 \(k\)。写出此杆的纵振动方程。

\begin{solution}
设纵向位移为 \(u(x,t)\)，杆的杨氏模量为 \(E\)，密度为 \(\rho\)，截面积为
\[
S(x)=kx^2.
\]
变截面杆纵振动的一般方程为
\[
\rho S(x)\frac{\partial^2u}{\partial t^2}
=\frac{\partial}{\partial x}
\left(E S(x)\frac{\partial u}{\partial x}\right).
\]
代入 \(S(x)=kx^2\)，并约去公共因子 \(k\)，得
\[
\rho x^2u_{tt}=E\frac{\partial}{\partial x}(x^2u_x)
=E(x^2u_{xx}+2xu_x).
\]
因此
\[
u_{tt}=a^2\left(u_{xx}+\frac{2}{x}u_x\right),
\qquad
a^2=\frac{E}{\rho},\quad 0<x<l.
\]
\end{solution}
\end{question}

\begin{question}
在铀块中，除了中子的扩散运动外，还存在中子的吸收和增殖过程。设在单位时间内、
单位体积中吸收和增殖的中子数均正比于该时刻、该处的中子浓度 \(u(\bm r,t)\)，
因而净增中子数可表示为 \(\alpha u(\bm r,t)\)，\(\alpha\) 为比例常数。
试导出 \(u(\bm r,t)\) 所满足的偏微分方程。
\begin{solution}
扩散通量满足 Fick 定律
\[
\bm J=-D\nabla u.
\]
守恒律给出
\[
\frac{\partial u}{\partial t}+\nabla\cdot\bm J=\alpha u.
\]
代入通量表达式，得到
\[
\frac{\partial u}{\partial t}
=D\nabla^2u+\alpha u.
\]
这就是带吸收、增殖源项的扩散方程；若 \(\alpha>0\) 表示净增殖，若 \(\alpha<0\) 表示净吸收。
\end{solution}
\end{question}

\begin{question}
由 Maxwell 方程组
\[
\begin{cases}
\nabla\cdot\bm E=\dfrac{\rho_e}{\varepsilon_0},\\
\nabla\times\bm E=-\dfrac{\partial\bm B}{\partial t},\\
\nabla\cdot\bm B=0,\\
\nabla\times\bm B=\varepsilon_0\mu_0\dfrac{\partial\bm E}{\partial t}+\mu_0\bm j,
\end{cases}
\]
以及恒等式
\[
\nabla^2\bm V\equiv
\nabla(\nabla\cdot\bm V)-\nabla\times(\nabla\times\bm V)
\]
推导电场强度 \(\bm E\) 和磁场强度 \(\bm B\) 所满足的微分方程，并说明由此可得出的结论。
\begin{solution}
对 Faraday 定律取旋度：
\[
\nabla\times(\nabla\times\bm E)
=-\frac{\partial}{\partial t}(\nabla\times\bm B).
\]
这一步的目的，是把一阶旋度方程变成含二阶空间导数的方程。由
\[
\nabla\times(\nabla\times\bm E)
=\nabla(\nabla\cdot\bm E)-\nabla^2\bm E
\]
以及
\[
\nabla\cdot\bm E=\frac{\rho_e}{\varepsilon_0},\qquad
\nabla\times\bm B=\varepsilon_0\mu_0\bm E_t+\mu_0\bm j,
\]
右端变为
\[
-\frac{\partial}{\partial t}(\nabla\times\bm B)
=-\varepsilon_0\mu_0\bm E_{tt}-\mu_0\bm j_t.
\]
因此
\[
\begin{split}
\nabla(\nabla\cdot\bm E)-\nabla^2\bm E
&=-\mu_0\frac{\partial\bm j}{\partial t}
-\varepsilon_0\mu_0\frac{\partial^2\bm E}{\partial t^2},\\
\nabla^2\bm E-\varepsilon_0\mu_0\frac{\partial^2\bm E}{\partial t^2}
&=\frac{1}{\varepsilon_0}\nabla\rho_e+\mu_0\frac{\partial\bm j}{\partial t}.
\end{split}
\]
对 Ampere-Maxwell 方程取旋度：
\[
\nabla\times(\nabla\times\bm B)
=\varepsilon_0\mu_0\frac{\partial}{\partial t}(\nabla\times\bm E)
+\mu_0\nabla\times\bm j.
\]
再用
\[
\nabla\times(\nabla\times\bm B)
=\nabla(\nabla\cdot\bm B)-\nabla^2\bm B=-\nabla^2\bm B,
\]
以及 \(\nabla\times\bm E=-\bm B_t\)，得到
\[
\nabla^2\bm B-\varepsilon_0\mu_0\frac{\partial^2\bm B}{\partial t^2}
=-\mu_0\nabla\times\bm j.
\]
在无源区域 \(\rho_e=0,\ \bm j=0\) 中，二者化为波动方程
\[
\nabla^2\bm E-\frac{1}{c^2}\bm E_{tt}=0,\qquad
\nabla^2\bm B-\frac{1}{c^2}\bm B_{tt}=0,
\qquad
c=\frac{1}{\sqrt{\varepsilon_0\mu_0}}.
\]
因此电磁场可以在真空中以速度 \(c\) 传播，这就是电磁波。
\end{solution}
\end{question}

\begin{question}
试求下述定解问题所描述的一维均匀细杆中平衡状态下的温度分布，以及整个细杆内单位时间所产生的热量、
杆的两端单位时间内分别散出的热量：
\[
\begin{cases}
\dfrac{\partial u}{\partial t}-\kappa\dfrac{\partial^2u}{\partial x^2}
=\kappa x^2,\qquad 0<x<L,\ t>0,\\
u(x,t)|_{x=0}=T,\qquad t\ge0,\\
\left.\dfrac{\partial u}{\partial x}\right|_{x=L}=0,\qquad t\ge0,\\
u(x,t)|_{t=0}=f(x),\qquad 0\le x\le L.
\end{cases}
\]
其中温度传导率 \(\kappa\) 以及 \(x=0\) 端的温度 \(T\) 都是常数。
设杆的密度为 \(\rho\)，比热容为 \(c\)。
\begin{solution}
平衡态与时间无关，记为 \(U(x)\)。由原方程得
\[
-\kappa U''(x)=\kappa x^2,
\]
即
\[
U''(x)=-x^2.
\]
两次积分：
\[
U'(x)=-\frac{x^3}{3}+C_1,\qquad
U(x)=-\frac{x^4}{12}+C_1x+C_2.
\]
边界条件给出
\[
U(0)=T,\qquad U'(L)=0.
\]
因此
\[
C_2=T,\qquad C_1=\frac{L^3}{3}.
\]
平衡温度分布为
\[
U(x)=T+\frac{L^3}{3}x-\frac{x^4}{12}.
\]
由热传导方程
\[
u_t=\kappa u_{xx}+\kappa x^2
\]
可知单位体积单位时间内的热源强度为
\[
\rho c\,\kappa x^2.
\]
故全杆单位时间产生的热量为
\[
Q_{\rm gen}=\int_0^L \rho c\kappa x^2\,dx
=\frac{\rho c\kappa L^3}{3}.
\]
热流密度为
\[
q=-\rho c\kappa U'(x).
\]
在 \(x=L\) 端，因 \(U'(L)=0\)，散热为
\[
Q_L=0.
\]
在 \(x=0\) 端，
\[
U'(0)=\frac{L^3}{3},
\]
热流指向负 \(x\) 方向，因此从 \(x=0\) 端散出的热量为
\[
Q_0=\rho c\kappa\frac{L^3}{3}.
\]
这与全杆内部产生的热量相等，符合稳态能量守恒。
\end{solution}
\end{question}

